\chapter{Deployment \& Production: Dari Laptop ke Cloud}
\label{chap:deployment}

Membuat skrip AI di Jupyter Notebook itu mudah. Tantangannya adalah membawanya ke tahap produksi agar bisa diakses oleh ribuan user. Bab ini membahas cara "membungkus" (containerize) dan men-deploy aplikasi AI.

\section{Membangun API dengan FastAPI}

Streamlit bagus untuk demo, tapi untuk backend produksi, **FastAPI** adalah standar industri karena cepat (asynchronous) dan otomatis membuat dokumentasi (Swagger UI).

\subsection{1. Struktur Project}
\begin{codeblock}
my-ai-app/
|-- main.py
|-- requirements.txt
|-- Dockerfile
`-- .env
\end{codeblock}

\subsection{2. Code: main.py}
\begin{codeblock}[language=Python]
from fastapi import FastAPI, HTTPException
from pydantic import BaseModel
from openai import OpenAI
import os
from dotenv import load_dotenv

load_dotenv()
client = OpenAI()

app = FastAPI(title="AI Service API")

# Definisi Schema Input
class ChatRequest(BaseModel):
    message: str
    model: str = "gpt-3.5-turbo"

# Definisi Schema Output
class ChatResponse(BaseModel):
    reply: str
    usage: dict

@app.get("/")
def read_root():
    return {"status": "AI Service is Online"}

@app.post("/chat", response_model=ChatResponse)
async def chat_endpoint(request: ChatRequest):
    try:
        response = client.chat.completions.create(
            model=request.model,
            messages=[{"role": "user", "content": request.message}]
        )
        return ChatResponse(
            reply=response.choices[0].message.content,
            usage=dict(response.usage)
        )
    except Exception as e:
        raise HTTPException(status_code=500, detail=str(e))

if __name__ == "__main__":
    import uvicorn
    uvicorn.run(app, host="0.0.0.0", port=8000)
\end{codeblock}

\subsection{3. Menjalankan Server}
\begin{codeblock}
pip install fastapi uvicorn openai python-dotenv
python main.py
\end{codeblock}
Buka \texttt{http://localhost:8000/docs} untuk melihat Swagger UI.

\section{Containerization dengan Docker}

Docker memastikan aplikasi Anda berjalan sama persis di laptop Anda dan di server cloud.

\subsection{1. Dockerfile}
Buat file bernama \texttt{Dockerfile} (tanpa ekstensi):

\begin{codeblock}
# Gunakan image Python ringan
FROM python:3.9-slim

# Set working directory
WORKDIR /app

# Salin requirements dulu (untuk caching layer)
COPY requirements.txt .
RUN pip install --no-cache-dir -r requirements.txt

# Salin sisa kode
COPY . .

# Expose port
EXPOSE 8000

# Perintah menjalankan aplikasi
CMD ["python", "main.py"]
\end{codeblock}

\subsection{2. Build \& Run}
\begin{codeblock}
# Build image
docker build -t my-ai-api .

# Run container
docker run -p 8000:8000 --env-file .env my-ai-api
\end{codeblock}

\section{Deployment ke Cloud}

\subsection{Opsi 1: Hugging Face Spaces (Gratis/Mudah)}
Cocok untuk demo Streamlit/Gradio.
\begin{enumerate}
    \item Buat akun di huggingface.co
    \item Create new Space $\rightarrow$ pilih SDK: Docker.
    \item Upload file Anda.
    \item HF akan otomatis build dan deploy.
\end{enumerate}

\subsection{Opsi 2: Railway / Render (PaaS)}
Cocok untuk API FastAPI.
\begin{enumerate}
    \item Push kode ke GitHub.
    \item Login ke Railway.app dengan GitHub.
    \item "New Project" $\rightarrow$ Pilih repo.
    \item Tambahkan Environment Variable (OPENAI\_API\_KEY).
    \item Deploy.
\end{enumerate}

\subsection{Opsi 3: AWS EC2 (Manual/Enterprise)}
\begin{enumerate}
    \item Sewa VPS (EC2 Instance).
    \item SSH ke server.
    \item Install Docker.
    \item \texttt{git pull} kode Anda.
    \item \texttt{docker run}.
    \item Setup Nginx dan Domain (untuk HTTPS).
\end{enumerate}

\section{Monitoring \& Observability}

Bagaimana kita tahu jika AI "berhalusinasi" di produksi?

\subsection{LangSmith (by LangChain)}
Platform untuk melacak setiap step (chain) dari aplikasi LLM.
\begin{codeblock}
export LANGCHAIN_TRACING_V2=true
export LANGCHAIN_API_KEY=ls__...
\end{codeblock}
Cukup set env var di atas, dan semua log LangChain akan muncul di dashboard LangSmith.

\subsection{Logging Sederhana}
Selalu log input dan output ke file atau database.

\begin{codeblock}[language=Python]
import logging

logging.basicConfig(filename='ai_logs.log', level=logging.INFO)

# Di dalam fungsi chat
logging.info(f"User: {request.message} | AI: {response_text}")
\end{codeblock}

\section{Security Checklist untuk Produksi}

\begin{itemize}
    \item [ ] **Jangan commit .env:** Gunakan \texttt{.gitignore}.
    \item [ ] **Rate Limiting:** Pasang pembatas request per IP agar tidak di-DDoS (bisa pakai library \texttt{slowapi} di FastAPI).
    \item [ ] **Input Validation:** Pastikan pesan user tidak terlalu panjang (misal max 1000 karakter) untuk mencegah serangan DoS biaya token.
    \item [ ] **API Key Rotation:** Ganti kunci API secara berkala.
\end{itemize}
